    \chapter{Polynomial Approximations to functions}
    With the help of a special identity $1-x^n = (1-x)(1+x+x^2+\cdots + x^{n-1})$, we can find a polynomial approximation to $-\ln(1-x)$ in the previous section. However, this may not help to find a polynomial approximation to other functions, which motivates us to find a more general method to find polynomial approximations to functions.
    \section{The Taylor Polynomial for a function}
    \begin{theoremenv}
        Let $f$ be a function with derivatives of order $n$ at the point $x=a$. Then there exists a unique polynomial $P$ of degree at most $n$ such that $P(a)=f(a), P'(a)=f'(a), \ldots, P^{(n)}(a) = f^{(n)}(a)$. Additionally, this polynomial is given by
        \[
        P(x) = \sum_{k=0}^n \frac{f^{(k)}(a)}{k!} (x-a)^k = f(a) + f'(a)(x-a) + \cdots + \frac{f^{(n)}(a)}{n!} (x-a)^n.
        \]
        and it is referred to as the \textbf{Taylor polynomial of degree $n$ for $f$ at $x=a$}.
    \end{theoremenv}
    \begin{notationenv}
        We often denote $P$ in the above theorem by $T_n(f)$ or $T_n(f;a)$ to emphasize the degree and the center of the Taylor polynomial.
    \end{notationenv}
    \begin{exampleenv}
        Here are some examples of Taylor polynomials:
        \begin{enumerate}[label=(\arabic*)]
            \item The Taylor polynomial of degree $n$ for $f(x)=e^x$ at $x=0$ is given by
            \[
            T_n(f) = 1 + x + \frac{x^2}{2!} + \cdots + \frac{x^n}{n!} = \sum_{k=0}^n \frac{x^k}{k!}.
            \]
            \item The Taylor polynomial of degree $2n+1$ for $f(x)=\sin x$ at $x=0$ is given by
            \[
            T_{2n+1}(f) = x - \frac{x^3}{3!} + \cdots + (-1)^n \frac{x^{2n+1}}{(2n+1)!} = \sum_{k=0}^n (-1)^k \frac{x^{2k+1}}{(2k+1)!}.
            \]
            \item The Taylor polynomial of degree $2n$ for $f(x)=\cos x$ at $x=0$ is given by
            \[
            T_{2n}(f) = 1 - \frac{x^2}{2!} + \cdots + (-1)^n \frac{x^{2n}}{(2n)!} = \sum_{k=0}^n (-1)^k \frac{x^{2k}}{(2k)!}.
            \]
        \end{enumerate}
    \end{exampleenv}
    \begin{theoremenv}
        Here are some properties of the Taylor polynomial:
        \begin{enumerate}[label=(\arabic*)]
            \item Linearity: if $c_1,c_2 \in \mathbb R$, then $$T_n(c_1 f + c_2 g) = c_1 T_n(f) + c_2 T_n(g).$$
            \item Differentiation: $$(T_n(f))' = T_{n-1}(f').$$
            \item Integration: if $\displaystyle F(x)=\int_a^x f(t)\, \dd t$, then \[T_{n+1}(F) = \int_a^x T_n(f(t))\, \dd t.\]
            \item Substitution: if $g(x)=f(cx)$ where $c\in \mathbb R$ is a constant, then \[T_n(g)(x;a) = T_n(f)(cx; ca).\]
            Particularly, if $a=0$, then $T_n(g)(x) = T_n(f)(cx)$.
        \end{enumerate}
    \end{theoremenv}
    \begin{exampleenv}
        With the help of the properties above, we continue to find the Taylor polynomial of degree $n$ for more functions at $x=0$:
        \begin{enumerate}[label=(\arabic*),start=4]
            \item The Taylor polynomial of degree $2n$ for $f(x)=\cosh x$ at $x=0$ is given by
            \[
            T_{2n}(f) = 1 + \frac{x^2}{2!} + \cdots + \frac{x^{2n}}{(2n)!} = \sum_{k=0}^n \frac{x^{2k}}{(2k)!}.
            \]
            \item The Taylor polynomial of degree $2n-1$ for $f(x)=\sinh x$ at $x=0$ is given by
            \[
            T_{2n-1}(f) = x + \frac{x^3}{3!} + \cdots + \frac{x^{2n-1}}{(2n-1)!} = \sum_{k=0}^{n-1} \frac{x^{2k+1}}{(2k+1)!}.
            \]
            \item Recall
            \[
            \frac1{1-x} = 1 + x + x^2 + \cdots + x^n + \frac{x^{n+1}}{1-x}:=P_n(x)+\frac{x^{n+1}}{1-x} \quad \text{for all } x\neq 1.
            \]
            Let $h(x) = x^n g(x)$ where $g(x) = \dfrac{x}{1-x}$, then we have
            \[
            f(x)=h(x) + P_n(x) \quad \text{for all } x\neq 1.
            \]
            Note that by checking the derivatives of $f$ and $P_n$ at $x=0$, we have
            \[
            f^{(k)}(0) = P_n^{(k)}(0) \quad \text{for all } k\in \{0,1,\ldots,n\}.
            \]
            Thus we have $T_n(x) = P_n(x)$. By $\displaystyle \int_0^x \frac{1}{1-t} \dd t = -\ln(1-x)$, we have
            \[
            T_{n+1}\left[ -\ln(1-x) \right] = \int_0^x T_n\left( \frac{1}{1-t} \right) \dd t = \int_0^x P_n(t) \dd t = \sum_{k=1}^{n+1} \frac{x^k}{k}.
            \]
        \end{enumerate}
    \end{exampleenv}
    \begin{theoremenv}
        Let $P_n$ be a polynomial of degree $n\ge 1$. Let $f$ and $g$ be two functions with derivatives of order $n$ at $x=0$ and assume that $$f(x)=P_n(x)+x^n g(x),$$ where $g(x)\to 0$ as $x\to 0$. Then $P_n$ is the Taylor polynomial $T_n(f)$ at 0.
    \end{theoremenv}
    \begin{remarkenv}
        the Taylor polynomial is the unique polynomial approximation whose error is \(x^ng(x)\) as \(x\to0\). Here, $x^n g(x)$ is called the \textbf{Peano form of the error term} of the Taylor polynomial $P_n$.
    \end{remarkenv}
    \begin{exampleenv}
        \begin{enumerate}[label=(\arabic*),start=7]
        \item We start again by
        \[
        \frac{1}{1-x} = 1 + x + x^2 + \cdots + x^n + \frac{x^{n+1}}{1-x}
        \]
        but replace $x$ by $(-x^2)$:
        \[
        \frac{1}{1+x^2} = 1 - x^2 + x^4 - \cdots + (-1)^n x^{2n} + \frac{(-1)^{n+1} x^{2n+1}}{1+x^2}\quad \forall x\in \R.
        \]
        Analogously, we let $P_{2n}(x) = 1 - x^2 + x^4 - \cdots + (-1)^n x^{2n}$ and $g(x) = \dfrac{(-1)^{n+1} x}{1+x^2}$, then we have
        \[
        \frac{1}{1+x^2} = P_{2n}(x) + x^{2n} g(x) \quad \text{for all } x\in \R.
        \]
        By the above theorem, we have $T_{2n}\left( \dfrac{1}{1+x^2} \right) = P_{2n}(x)$. By $\displaystyle \int_0^x \frac{1}{1+t^2} \dd t = \arctan x$, we have
        \[
        T_{2n+1}(\arctan x) = \int_0^x T_{2n}\left( \frac{1}{1+t^2} \right) \dd t = \int_0^x P_{2n}(t) \dd t = \sum_{k=0}^n \frac{(-1)^k x^{2k+1}}{2k+1}.
        \]
        \end{enumerate}
    \end{exampleenv}
    \begin{definitionenv}
        We call the function $E_n(x) := f(x) - T_n(f(x))$ the \textbf{error term} or the \textbf{remainder} of the Taylor polynomial $T_n(f)$.
    \end{definitionenv}
    \begin{exampleenv}
        \begin{enumerate}[label=(\arabic*),start=8]
            \item Suppose $f$ has continuous second derivative $f''$ in some neighborhood of $a$. Then in this neighborhood of $a$, we know:
            \[
            \begin{aligned}
                E_1(x) &= f(x) - T_1(f) \\
                &= f(x) - \left[ f(a) + f'(a)(x-a) \right] \\
                &= \int_a^x (f'(t) - f'(a))\, \dd t\quad (\text{IBP with } u=f'(t)-f'(a),v=t-x ) \\
                &= \left[ f'(t) - f'(a) \right](t-x)\Big|_{t=a}^x - \int_a^x f''(t)(t-x)\, \dd t\\
                &= \int_a^x f''(t)(x-t)\, \dd t.
            \end{aligned}
            \]
            If, additionally, we know $m\le f''(t) \le M$ for any $t$ in the considered neighborhood of $a$, then we have
            \[
            \frac{m(x-a)^2}{2} \le E_1(x) = \int_a^x f''(t)(x-t)\, \dd t \le \frac{M(x-a)^2}{2}.
            \]
            By the squeeze theorem, we have $\displaystyle \lim_{x\to a} E_1(x) = 0$. This example motivates us to find a more general formula for the error term $E_n$ as a integral.
        \end{enumerate}
    \end{exampleenv}
    \begin{theoremenv}\label{Taylor formula with error term}
        Assume $f$ has a continuous derivative of order $n+1$ in some interval containing $a$. Then for every $x$ in this interval, we have the Taylor formula:
        \[
        f(x) = T_n(f) + E_n(x) = \sum_{k=0}^n \frac{f^{(k)}(a)}{k!} (x-a)^k + \frac{1}{n!} \int_a^x f^{(n+1)}(t) (x-t)^n\, \dd t.
        \]
        If, additionally, there exists $m$ and $M$ such that $m \le f^{(n+1)}(t) \le M$ for any $t$ in the considered interval, then we have the following bounds/estimates for the error term:
        \[
        \frac{m(x-a)^{n+1}}{(n+1)!} \le E_n(x) = \frac{1}{n!} \int_a^x f^{(n+1)}(t) (x-t)^n\, \dd t \le \frac{M(x-a)^{n+1}}{(n+1)!} \ (\forall x>a).
        \]
        and
        \[
        \frac{m(a-x)^{n+1}}{(n+1)!} \le (-1)^{n+1} E_n(x)  \le \frac{M(a-x)^{n+1}}{(n+1)!} \quad \forall x<a.
        \]
        \end{theoremenv}
        \begin{remarkenv}
            Here, $\displaystyle E_n(x) = \frac{1}{n!} \int_a^x f^{(n+1)}(t) (x-t)^n\, \dd t$ is called the \textbf{integral form of the error term} of the Taylor polynomial $T_n(f)$.
        \end{remarkenv}
        \begin{exampleenv}
            We go back to the example of $\displaystyle f(x) = e^x=\sum_{k=0}^n \frac{x^k}{k!} + E_n(x)$ and $a=0$. Note for any $x\in [b,c]$, we know $e^b \le e^x \le e^c$. Particularly, if $b=0$ and $c=x=1$, then we have
            \[
            \frac{1}{(n+1)!} \le E_n(1) \le e\cdot \frac{1}{(n+1)!} < \frac{3}{(n+1)!}.
            \]
            Thus we have $\displaystyle e=\sum_{k=0}^n \frac{1}{k!} + E_n(1)$.
        \end{exampleenv}
        \begin{remarkenv}
            Recall the error term
            \[
            E_n(x) = \frac{1}{n!} \int_a^x f^{(n+1)}(t) (x-t)^n\, \dd t.
            \]
            Note that for any $t\in [a,x]$, the term $(x-t)^n$ never changes sign. Back in Theorem~\ref{Taylor formula with error term}, we require $f^{(n+1)}$ to be continuous on the considered interval. By the weighted Mean Value Theorem for integrals, we can find some $c\in [a,x]$ such that
            \[
            \begin{aligned}
                E_n(x) &= \frac{1}{n!} \int_a^x f^{(n+1)}(t) (x-t)^n\, \dd t\\
                &= \frac{1}{n!} f^{(n+1)}(c) \int_a^x (x-t)^n\, \dd t\\
                &= \frac{1}{(n+1)!} f^{(n+1)}(c) (x-a)^{n+1}\quad \text{for some } c\in [a,x].
            \end{aligned}
            \]
            This is called the \textbf{Lagrange form of the error term}.

            Suppose $f^{(n+1)}$ is continous on some interval $[a-r,a+r]$ where $r>0$, then $f^{(n+1)}$ is bounded on $[a-r,a+r]$ and there exists $M>0$ such that $|f^{(n+1)}(t)| \le M$ for any $t\in [a-r,a+r]$. Then
            \[
            |E_n(x)| \le M \cdot \frac{|x-a|^{n+1}}{(n+1)!} \quad \text{for any } x\in [a-r,a+r].
            \]
            For $x\neq a$, we know $$0 \le \big|\frac{E_n(x)}{(x-a)^{n}}\big| \le \frac{M}{(n+1)!}|x-a|.$$ By the squeeze theorem, we know $\displaystyle \frac{E_n(x)}{(x-a)^{n}} \to 0$ as $x \to a$. We say that the error term $E_n$ is of small order that $(x-a)^n$ as $x\to a$, and we often write $$E_n(x) = o\left( (x-a)^n \right) \quad \text{as } x\to a.$$
            Here, $E_n(x)$ is called the \textbf{Peano form of the error term} of the Taylor polynomial $T_n(f)$.
        \end{remarkenv}
        \begin{definitionenv}
            Assume $g(x)\neq 0$ for $x\neq a$ in some interval containing $a$. The notation $f(x) = o(g(x))$ as $x\to a$ means that $\displaystyle \lim_{x\to a} \frac{f(x)}{g(x)} = 0$.
        \end{definitionenv}
        \begin{exampleenv}
            We can rewrite Taylor's formula in o-notation as
            \[
            f(x) = T_n(f) + o\left( (x-a)^n \right) \quad \text{as } x\to a.
            \]
        \end{exampleenv}
        \begin{theoremenv}
            We deduce some useful properties of the o-notation: as $x\to a$,
            \begin{enumerate}[label=(\arabic*)]
                \item $o(g(x)) \pm o(g(x)) = o(g(x)).$
                \item $o(cg(x)) = o(g(x))$ for any real constant $c\neq 0$.
                \item $f(x) \cdot o(g(x)) = o(f(x)g(x))$ for $f(x)\neq 0$.
                \item $o(o(g(x))) = o(g(x))$.
                \item $\displaystyle \frac{1}{1+g(x)} = 1 - g(x) + o(g(x))$ if $g(x)\to 0$ as $x\to a$.
            \end{enumerate}
        \end{theoremenv}
        \begin{proofenv}
            We only prove (1) and (5) here since the others are analogous.
            \begin{enumerate}[label=(\arabic*)]
                \item Let $f_1(x) = o(g(x))$ and $f_2(x) = o(g(x))$, then $\displaystyle \lim_{x\to a} \frac{f_1(x)}{g(x)} = 0$ and $\displaystyle \lim_{x\to a} \frac{f_2(x)}{g(x)} = 0$. Thus we have
                \[
                \lim_{x\to a} \frac{f_1(x) \pm f_2(x)}{g(x)} = \lim_{x\to a} \frac{f_1(x)}{g(x)} \pm \lim_{x\to a} \frac{f_2(x)}{g(x)} = 0.
                \]
                Thus $o(g(x)) + o(g(x)) = f_1(x) + f_2(x) = o(g(x))$.
                \item[(5)] Recall
                \[
                \frac{1}{1+g(x)} = 1 - g(x) + \frac{g(x)^2}{1+g(x)}\quad \text{for any } g(x)\neq -1.
                \]
                Thus we have
                \[
                \lim_{x\to a} \frac{\frac{1}{1+g(x)} - (1 - g(x))}{g(x)} = \lim_{x\to a} \frac{g(x)^2}{(1+g(x))g(x)} = \lim_{x\to a} \frac{g(x)}{1+g(x)} = 0.
                \]
                Thus we have $\displaystyle \frac{1}{1+g(x)} = 1 - g(x) + o(g(x))$.
            \end{enumerate}
        \end{proofenv}
        \begin{exampleenv}
            We know 
            \[
            \cos x = 1 - \frac{x^2}{2!} + \cdots + \frac{(-1)^n x^{2n}}{(2n)!} + o(x^{2n+1}) \quad \text{as } x\to 0.
            \]
                Let $g(x) = \cos x -1$, then $g(x)\to 0$ as $x\to 0$. By the properties of o-notation, we have
                \[
                \frac{1}{\cos x} = \frac{1}{1 + g(x)} = 1 - g(x) + o(g(x)) = 1 +\frac{x^2}{2} + o(x^2) + o(x^3) \quad \text{as } x\to 0.
                \]
            Note $o(x^2)+o(x^3)=o(x^2)$, so we have
            \[
            \tan x = \frac{\sin x}{\cos x} = \left( x - \frac{x^3}{3!} + o(x^4) \right) \left( 1 + \frac{x^2}{2} + o(x^3) \right) = x + \frac{x^3}{3} + o(x^3) \quad \text{as } x\to 0.
            \]
        \end{exampleenv}
        \begin{exampleenv}
            With Taylor's formula, we can calculate the following limit:
            \begin{enumerate}
                \item \(\displaystyle
            \lim_{x\to 0} \frac{a^x - b^x}{x} \quad \text{where } a,b>0.
            \)

            Note that
            \[
            e^x = 1 + x + o(x) \quad \text{as } x\to 0.
            \]
            Thus we have
            \[
            a^x - b^x = e^{x\ln a} - e^{x\ln b} = x(\ln a - \ln b) + o(x) \quad \text{as } x\to 0.
            \]
            By definition of o-notation, we have
            \[
            \lim_{x\to 0} \frac{a^x - b^x}{x} = \ln a - \ln b.
            \]
                \item \(\displaystyle
            \lim_{x\to 0} \frac{\ln(1+ax)}{x} \quad \text{where } a\in \mathbb{R}.
            \)

            If $a\neq 0$, note that
            \[
            \ln(1+ax) = ax + o(x) \quad \text{as } x\to 0.
            \]
            Thus we have
            \[
            \frac{\ln(1+ax)}{x} = a + o(1) \quad \text{as } x\to 0 \implies \lim_{x\to 0} \frac{\ln(1+ax)}{x} = a.
            \]
            But, what if $a=0$? 
            \end{enumerate}
        \end{exampleenv}

        \section{Indeterminate forms and L'Hôpital's Rule}
        \begin{theoremenv}
            Assume $f$ and $g$ have derivatives $f'(x)$ and $g'(x)$ for every $x$ in an open interval $(a,b)$ and suppose that $\displaystyle \lim_{x\to a^+} f(x) = \lim_{x\to a^+} g(x) = 0$. If $g'(x) \neq 0$ for every $x\in (a,b)$ and the limit $\displaystyle \lim_{x\to a^+} \frac{f'(x)}{g'(x)}$ exists and have value $L$, then we have $$\lim_{x\to a^+} \frac{f(x)}{g(x)} = \frac{f'(x)}{g'(x)} = L.$$
        \end{theoremenv}
        \begin{proofenv}
            Note that L'ôpital's Rule does require $f,g,f',g'$ to be well-defined for $x\in (a,b)$. So we shall ``extend'' the definition of $f$ and $g$ ``a bit'' by letting
            \[
            F(x) = \begin{cases}
            f(x) & \text{if } x\in (a,b)\\
            0 & \text{if } x=a
            \end{cases}
            \quad \text{and} \quad
            G(x) = \begin{cases}
            g(x) & \text{if } x\in (a,b)\\
            0 & \text{if } x=a
            \end{cases}
            \]
            We can apply Cauchy's Mean Value Theorem to the interval $[a,x]$ for $F$ and $G$ to get
            \[
            \frac{F(x) - F(a)}{G(x) - G(a)} = \frac{F'(c)}{G'(c)} \quad \text{for some } c\in (a,x).
            \]
            Simply note that $F(a) = G(a) = 0$, we have
            \[
            \frac{f(x)}{g(x)} = \frac{F(x)}{G(x)} = \frac{F'(c)}{G'(c)} = \frac{f'(c)}{g'(c)} \quad \text{for some } c\in (a,x).
            \]
            Note $g(x) \neq 0$ for $x\in (a,b)$ otherwise $G(x) = 0 = G(a)$ and by Role's Theorem, there exists some $d\in (a,x)$ such that $G'(d) = 0 = g'(d)$, which is a contradiction. Thus as $x\to a^+$, we have $c\to a^+$ and $$\lim_{x\to a^+} \frac{f(x)}{g(x)} = \frac{f'(x)}{g'(x)} = L.$$
        \end{proofenv}
        \begin{remarkenv}
            The "left-sided limit" and the "two-sided limit" versions of L'Hôpital's Rule also hold provided that the derivatives are defined in the needed open interval.
            
            More versions of L'Hôpital's Rule can be derived for indeterminate forms of the type $\infty/\infty$, $0\cdot \infty$, $\infty - \infty$, $0^0$, $1^\infty$ and $\infty^0$ or for ``limit notation'' like $\displaystyle \lim_{x\to \infty } \frac{f(x)}{g(x)}$, $\displaystyle \lim_{x\to a} \frac{f(x)}{g(x)} = \infty$ or $\displaystyle \lim_{x\to \infty} \frac{f(x)}{g(x)} = -\infty$.
        \end{remarkenv}
        \begin{exampleenv}
            We use L'Hôpital's Rule to calculate some previous limits:
            \begin{enumerate}[label=(\arabic*)]
                \item Calculate \(\displaystyle
            \lim_{x\to 0} \frac{a^x - b^x}{x}\text{ where } a,b>0.
            \) Note that
            \[
            \lim_{x\to 0} a^x = \lim_{x\to 0} b^x = 1.
            \]
            Thus we have
            \[
            \lim_{x\to 0} \frac{a^x - b^x}{x} = \lim_{x\to 0} \frac{a^x \ln a - b^x \ln b}{1} = \ln a - \ln b.
            \]
                \item Calculate \(\displaystyle
            \lim_{x\to 0} \frac{\ln(1+ax)}{x} \text{ where } a\in \mathbb{R}.
            \) We talk about the case $a=0$ here. Note that
            \[
            \lim_{x\to 0} \ln(1+ax) = \ln 1 = 0.
            \]
            Thus we have
            \[
            \lim_{x\to 0} \frac{\ln(1+ax)}{x} = \lim_{x\to 0} \frac{a}{1+ax} = a.
            \]
            \item $\displaystyle \lim_{x\to 0} \frac{x-\tan x}{x-\sin x} = \lim_{x\to 0}\frac{1-\sec^2 x}{1-\cos x} =-\lim_{x\to 0} \frac{\cos x+1}{\cos^2 x}= -2$.
            \end{enumerate}
        \end{exampleenv}
        \begin{definitionenv}
            We introduce extended definitions of limits and ``limits with infinity'' here:
            \begin{itemize}
                \item The symbol
                \[
                \lim_{x\to \infty} f(x) = L \quad \text{or} \quad f(x) \to L \quad \text{as } x\to \infty
                \]
                means that for every $\varepsilon > 0$, there exists some $M>0$ such that \[|f(x) - L| < \varepsilon \quad \text{for any } x > M. \]
                \item The symbol $\displaystyle \lim_{x\to a} f(x) = \infty$ means that for every $M>0$, there exists some $\delta > 0$ such that $f(x) > M$ whenever $0 < |x-a| < \delta$. Similarly, the symbol $\displaystyle \lim_{x\to a} f(x) = -\infty, \displaystyle \lim_{x\to a^+} f(x) = \pm \infty$ and $\displaystyle \lim_{x\to a^-} f(x) = \pm \infty$ can be defined analogously.
                \item The symbol $\displaystyle \lim_{x\to \infty} f(x) = \infty$ means that for every $M>0$, there exists some $N>0$ such that $f(x) > M$ whenever $x > N$. Similarly, the symbol $\displaystyle \lim_{x\to \infty} f(x) = -\infty$ can be defined analogously.
            \end{itemize}
        \end{definitionenv}
        \begin{remarkenv}
            Note, if $g(t) = f(1/t)$ and $\displaystyle \lim_{t\to 0^+} g(t) = A$, then $\displaystyle \lim_{x\to \infty} f(x) = A$. Similarly, if $\displaystyle \lim_{t\to 0^-} g(t) = B$, then $\displaystyle \lim_{x\to -\infty} f(x) = B$.
        \end{remarkenv}
        \begin{remarkenv}
            With the extended definitions of limits, we can also apply L'Hôpital's Rule to calculate limits like $\displaystyle \lim_{x\to \infty} \frac{f(x)}{g(x)}$ provided that the derivatives are defined in the needed open interval and the limit $\displaystyle \lim_{x\to \infty} \frac{f'(x)}{g'(x)}$ exists.
        \end{remarkenv}
        \begin{theoremenv}
            Suppose $a,b>0$ then $\displaystyle \lim_{x\to +\infty} \frac{(\ln x)^b}{x^a} = 0$ and $\displaystyle \lim_{x\to +\infty} \frac{x^b}{e^{ax}} = 0$.
        \end{theoremenv}
        \begin{proofenv}
            Suppose $c>0$ and $t\ge 1$. Then $\dfrac{1}{t} \le \dfrac{1}{t^{1-c}}$. For any $x>1$, we have
            \[
            0<\ln x = \int_1^x \frac{1}{t} \dd t \le \int_1^x \frac{1}{t^{1-c}} \dd t = \frac{x^c - 1}{c} < \frac{x^c}{c}.
            \]
            Thus we have
            \[
            0 < \frac{(\ln x)^b}{x^a} < \frac{x^{cb-a}}{c^b} \quad \text{for any } x>1.
            \]
            Since $c$ is arbitrary, we choose $c=\dfrac{a}{2b}$ to get $x^{cb-a} = x^{-\frac{a}{2}} \to 0$ as $x\to +\infty$. By the squeeze theorem, we have $\displaystyle \lim_{x\to +\infty} \frac{(\ln x)^b}{x^a} = 0$.

            For the second limit, we let $t=e^x$ and we have
            \[
            \lim_{x\to +\infty} \frac{x^b}{e^{ax}} = \lim_{t\to +\infty} \frac{(\ln t)^b}{t^a} = 0.
            \]
        \end{proofenv}
    \begin{exampleenv}
        \[
        \lim_{x\to 0} x^x = \lim_{x\to 0} e^{x\ln x} = e^{\lim_{x\to 0} x\ln x} = e^0 = 1.
        \]
    \end{exampleenv}
    \begin{exampleenv}
        \[
        \lim_{x\to \infty} x^{1/x} = \lim_{x\to 0} \frac{1}{x^x} = \frac{1}{\displaystyle \lim_{x\to 0} x^x} = \frac{1}{1} = 1.
        \]
    \end{exampleenv}
    \newpage